% =====================================================================
%  MODELLO: VERBALE ESTERNO  (incontro con un proponente o con i docenti)
%
%  1. Copiate questo file nella cartella della milestone in corso:
%       Candidatura/verbali/esterni/   RTB/verbali/esterni/   PB/verbali/esterni/
%     e chiamatelo VE-AAAA-MM-GG.tex (data dell'incontro), es.
%       RTB/verbali/esterni/VE-2026-11-05.tex
%     Dopo il merge la CI mette accanto il PDF: VE-2026-11-05.pdf
%     La copia firmata dal proponente salvatela accanto con un altro
%     nome: VE-2026-11-05-firmato.pdf (vedi GUIDA.md, ricetta 3)
%  2. Cambiate i campi qui sotto (dopo il simbolo % c'è un commento).
%  3. Scrivete il contenuto dopo \begin{document}.
%  Rispetto al verbale interno ci sono in più: i partecipanti esterni,
%  domande e risposte e, in fondo, il riquadro per la firma.
%  Guida completa: GUIDA.md (ricetta 2).
% =====================================================================
\documentclass[verbale=esterno]{gpt}

\titolo{Incontro con l'azienda}    % es. Chiarimenti sul capitolato C2
\data{2026-10-22}                  % data dell'incontro: AAAA-MM-GG
\orario{16:00}{17:00}              % ora di inizio, ora di fine
\luogo{Microsoft Teams}            % oppure: sede dell'azienda, Discord...
\redattore{corro}                  % chi scrive il verbale (la sua chiave)
\verificatore{cescobolla}          % chi lo controlla (diverso dal redattore)

% Chi faceva parte del gruppo il giorno dell'incontro.
% Lasciate così, a meno che il gruppo sia cambiato.
\componenti{biasini,cescobolla,corro,guglielmo,mazzurana,spagnoli}
\assenti{}                         % chi mancava, es. \assenti{spagnoli}

% Persone esterne presenti: una riga per persona
\esterno{Nome Cognome}{Ruolo, Azienda}
% \esterno{Altro Nome}{Altro ruolo, Azienda}

% Chi firma il verbale per l'azienda (compare nel riquadro firma finale)
\firma{Nome Cognome}{Azienda}

% Destinatari: ai docenti aggiungete l'azienda (\newline = a capo)
\destinatari{Prof.~Tullio Vardanega\newline Prof.~Riccardo Cardin\newline Azienda}

% Solo DOPO l'assegnazione del capitolato (togliete il % davanti):
% \capitolato{Cn -- Nome del capitolato (Azienda)}

% ---------------------------------------------------------------------
%  REGISTRO DELLE MODIFICHE  --  la riga più recente va IN CIMA.
%  La versione del documento è quella della PRIMA riga.
%
%  \modifica{versione}{data}{autore}{verificatore}{cosa è cambiato}
%
%  Quando il Verificatore approva, aggiunge lui IN CIMA questa riga:
%  \modifica{1.0.0}{2026-10-23}{corro}{cescobolla}{Verifica e approvazione}
% ---------------------------------------------------------------------
\modifica{0.1.0}{2026-10-22}{corro}{--}{Prima stesura}

\begin{document}

\section{Ordine del giorno}
\begin{enumerate}
  \item Presentazione del gruppo.
  \item Chiarimenti sui requisiti.
\end{enumerate}

\section{Domande e risposte}
% \domanda{la domanda fatta}{la risposta ricevuta}
\domanda{Prima domanda?}{Risposta ricevuta.}
\domanda{Seconda domanda?}{Risposta ricevuta.}

\section{Decisioni}
\begin{decisioni}
  \decisione{Testo della decisione presa.}
\end{decisioni}

\section{Azioni da intraprendere}
\begin{azioni}
  % \azione[numero issue]{cosa fare}{chi (chiavi, separate da virgole)}{entro quando}
  \azione{Descrizione dell'attività.}{tutti}{30/10/2026}
\end{azioni}

\end{document}
